Fix \(\mathbf M\in\mathcal C(\mathcal P)\). We verify the transition invariant in phase order.

For \(\mathsf{NORM}\), the attached derivation is a composition of consistency, unnesting, contradiction removal, Boolean disjoint-normal-form identities, and finite disjoint additivity. Each constructor preserves the event probability, so replacing the node preserves every root that reaches it. For \(\mathsf{ANCESTOR}\), the omitted variables range over their complete finite declared alphabets, and therefore
\[
P(e)=\sum_xP(e,x),
\]
which is the checked ancestor marginalization identity. For \(\mathsf{SPLIT}\), the compiler enumerates all assignments and the districts of \(G[H]\); \(\mathrm{ass:district\text{-}factorization}\) identifies the ancestor-closed probability with the checked finite sum of products. Hence the first three transitions preserve both root denotations.

At \(\mathsf{EXTRACT}\), a retained raw row has a checked proof both of district locality and of \(s=\star\) or \(D\cap\Delta_s=\varnothing\). Typed-menu realizability and finite recursive evaluation give
\[
P_a^s(v)=r^{\mathrm{raw}}_{s,a,D,v}q_D^s.
\]
The admissibility proof and district-kernel invariance replace \(q_D^s\) by \(q_D^\star\). For every retained derived row, \(\mathrm{lem:district\text{-}cell\text{-}linearity}\) gives
\[
K_{s,a,D}(c\mid h)=r_{s,a,D,h,c}q_D^\star.
\]
Together with \(\mathbf 1^\top q_D^\star=1\), stacking all accepted equalities yields
\[
A_Dq_D^\star(\mathbf M)=m_D.
\]
Both the row-needed node and its face-needed successor denote \(bq_D^\star(\mathbf M)\), while rejected candidates create no asserted equality. Thus extraction preserves denotation.

At \(\mathsf{FACE\_ID}\), the exact certificate supplies
\[
b=\lambda^\top A+\sum_{\omega\in Z}\zeta_\omega e_\omega,
\qquad Aq_D^\star=m,
\qquad q_D^\star(\omega)=0\quad(\omega\in Z).
\]
Consequently
\[
bq_D^\star=\lambda^\top m,
\]
which is exactly the resolved-node denotation. This calculation uses only inclusion of the global projection in the row face, not surjectivity onto that face. The \(\mathsf{MARGINALIZE}\) transition preserves denotation by finite additivity, and \(\mathsf{PRODUCT}\) preserves it by district factorization. This proves pointwise preservation for all seven nonterminal transformation rules.

The \(\mathsf{CONDITION}\) transition records the same resolved numerator and denominator. It constructs a quotient only when the independent certificate \(\pi\) proves strict positivity throughout \(\mathcal C(\mathcal P)\); the zero and failure branches perform no division.

For the operational assertions, the first seven rules respectively decrease
\[
r_N,r_A,r_S,r_X,r_F,r_{\Sigma},r_{\Pi}
\]
while fixing all earlier coordinates of \(\operatorname{rank}_{\mathrm{lex}}\), and no rule creates an earlier-phase node. Condition decreases \(r_C\) after the other coordinates vanish. Fixed priority, canonical finite orders, and the least-applicable-node convention therefore select at most one successor. Conversely, a well-typed nonterminal configuration has a least unresolved phase; its rule either applies or invokes the specified typed failure constructor for the failed partial operation. Hence \(\operatorname{STUCK}\) is unreachable.